\documentclass[11pt,a4paper]{article}
\usepackage[margin=20mm]{geometry}
\usepackage{amsmath,amssymb,booktabs}
\input{glyphtounicode}
\pdfgentounicode=1
\usepackage[hidelinks]{hyperref}
\setlength{\parindent}{0pt}
\setlength{\parskip}{5pt}
\setlength{\abovedisplayskip}{5pt}
\setlength{\belowdisplayskip}{5pt}
\newcommand{\der}[1]{\frac{\mathrm d #1}{\mathrm d\alpha}}
\begin{document}
\begin{center}
{\Large Equation 15: the concrete substitution map}
\end{center}

The upper line has already combined the lower cases. There is no remaining
$\mathrm dx_\ell/\mathrm d\alpha$ slot in it. To see exactly how they connect,
take two uncapped lanes, one capped lane, and an optional fourth lane
that remains at zero:
\[
 \mathcal A=\{1,2\},\qquad \mathcal U=\{3\},\qquad x_4=0.
\]

\textbf{First map the categories to the upper formula.}
\begin{center}
\renewcommand{\arraystretch}{1.5}
\begin{tabular}{@{}lll@{}}
\toprule
Upper-formula slot & Replacement & Reason\\
\midrule
$|\mathcal A|$ & $2$ & Lanes 1 and 2 are uncapped.\\
$\displaystyle\sum_{\mathcal U}\der{u_\ell}$
 & $\displaystyle\der{u_3}$ & Only lane 3 is capped.\\
$\displaystyle\sum_{\mathcal A}\der{h_\ell}$
 & $\displaystyle\der{h_1}+\der{h_2}$ & Only lanes 1 and 2 enter this sum.\\
\bottomrule
\end{tabular}
\end{center}
Thus the upper line becomes
\[
 \boxed{\der{\lambda}
 =\frac12\left(\der{V}-\der{u_3}+\der{h_1}+\der{h_2}\right).}
\]

\textbf{Here is precisely where the lower cases were substituted.}
The starting point is the derivative of $x_1+x_2+x_3+x_4=V$:
\begin{align*}
 \der{V}
 &=\der{x_1}+\der{x_2}+\der{x_3}+\der{x_4}\\
 &=\underbrace{\left(\der{\lambda}-\der{h_1}\right)}_{\text{lane 1: uncapped}}
  +\underbrace{\left(\der{\lambda}-\der{h_2}\right)}_{\text{lane 2: uncapped}}
  +\underbrace{\der{u_3}}_{\text{lane 3: capped}}
  +\underbrace{0}_{\text{lane 4: zero}}\\
 &=2\der{\lambda}-\der{h_1}-\der{h_2}+\der{u_3}.
\end{align*}
Rearranging this last line gives the boxed upper formula.

\textbf{Now calculate: upper first, then lower.}
Suppose the known input sensitivities are
\[
 \der{V}=0.3,\quad \der{u_3}=0.1,\quad
 \der{h_1}=0.2,\quad \der{h_2}=0.
\]
The upper formula gives
$\mathrm d\lambda/\mathrm d\alpha=(0.3-0.1+0.2+0)/2=0.2$.
Insert this result into each lower case:
\[
 \underbrace{\der{x_1}=0.2-0.2=0}_{\text{uncapped: level minus load}},\quad
 \underbrace{\der{x_2}=0.2-0=0.2}_{\text{uncapped: level minus load}},\quad
 \underbrace{\der{x_3}=0.1}_{\text{follows its cap}},\quad
 \underbrace{\der{x_4}=0}_{\text{stays zero}}.
\]
Check: the four individual gradients sum to $0+0.2+0.1+0=0.3$, exactly
the given total-admission gradient. These are sensitivities, not mass
values; the example assumes the lane categories remain unchanged.
\end{document}
